\section*{Chapter \ref{ch:g7:geological-time} --- The Ages of the Earth}

\begin{solution}{exo:g7:geological-time:1}
(a) \omterm{def:g7:geological-time:unconformity}{unconformity}; (b) inclusion; (c) superposition; (d) \omterm{def:g7:sediments-and-fossils:fossil}{fossil}
succession; (e) cross-cutting relationships.
\end{solution}

\begin{solution}{exo:g7:geological-time:2}
(a) The red sandstone, at the bottom, is the oldest; the yellow
sandstone, at the top, is the youngest. (b) The \omterm{def:g7:geological-time:superposition}{principle of
superposition}. (c) The layers have not been turned over since they
formed.
\end{solution}

\begin{solution}{exo:g7:geological-time:3}
(a) Paleozoic, Mesozoic, Cenozoic. (b) The Cenozoic. (c) The Mesozoic.
(d) The Jurassic comes right after the Triassic; the Cretaceous comes
right before the Paleogene.
\end{solution}

\begin{solution}{exo:g7:geological-time:4}
(a) False: \omterm{def:g7:geological-time:relative-dating}{relative dating} only puts \omterm{def:g3:stones-pebbles-sand:rock}{rocks} in order; an age in years is
an \omterm{def:g7:geological-time:relative-dating}{absolute age}. (b) True. (c) False: the \omterm{def:g3:stones-pebbles-sand:pebble}{pebbles} are older than the
conglomerate that holds them (\omterm{def:g7:geological-time:inclusion}{principle of inclusion}). (d) True
(\omterm{def:g7:geological-time:fossil-succession}{principle of fossil succession}).
\end{solution}

\begin{solution}{exo:g7:geological-time:5}
(a) Younger: it cuts layer 3, and a \omterm{def:g3:stones-pebbles-sand:rock}{rock} that cuts another is younger
(cross-cutting). (b) Older: layer 4 lies over it without being cut.
(c) Layer 1, layer 2, layer 3, the granite, layer 4. (d) The \omterm{def:g3:stones-pebbles-sand:rock}{rocks} above
it were worn away by \omterm{def:g5:water-shapes-the-land:erosion}{erosion} until the granite reached the surface;
layer 4 was then laid down on this \omterm{def:g5:water-shapes-the-land:erosion}{erosion} surface, an \omterm{def:g7:geological-time:unconformity}{unconformity}.
\end{solution}

\begin{solution}{exo:g7:geological-time:6}
(a) $538.8 - 251.9 = 286.9$ million years. (b) $251.9 - 66 = 185.9$
million years. (c) 66 million years. (d) The Paleozoic is the longest,
the Cenozoic the shortest (so far).
\end{solution}

\begin{solution}{exo:g7:geological-time:7}
(a) $4540 - 538.8 = 4001.2$ million years. (b) $4001.2 \div 4540
\approx 0.881$: about \textbf{88~\%}. (c) Before the Paleozoic, animals
with shells were not yet common, and soft living things almost never
become \omterm{def:g7:sediments-and-fossils:fossil}{fossils}; shells and bones, which fossilise well, are common only
from the Paleozoic on.
\end{solution}

\begin{solution}{exo:g7:geological-time:8}
(a) $4540 \div 100 = 45.4$ million years per metre. (b) $538.8 \div 45.4
\approx 11.9$: about \qty{11.9}{m} before the finish line. (c) $66 \div
45.4 \approx 1.45$: about \qty{1.45}{m}. (d) $0.3 \div 45.4 \approx
0.0066$~m, that is about \qty{0.66}{cm}, less than one centimetre: a
runner would hardly notice it.
\end{solution}

\begin{solution}{exo:g7:geological-time:9}
(a) The Silurian came first. Yes: the older \omterm{def:g3:stones-pebbles-sand:rock}{rock} lies below the younger
one. (b) The grey layers are laid down flat; they are pushed up on end;
they are worn away; the red \omterm{def:g3:stones-pebbles-sand:sand}{sand} is laid down on top. (c) It is a
surface of \omterm{def:g5:water-shapes-the-land:erosion}{erosion} between \omterm{def:g3:stones-pebbles-sand:rock}{rocks} of very different ages: the long time
during which the grey layers were raised and worn away left no layer
behind.
\end{solution}

\begin{solution}{exo:g7:geological-time:10}
(a) The two shales with ammonites are of the same age, and so are the
two limestones with \omterm{def:g2:water-on-earth:ocean}{sea} urchins: same \omterm{def:g7:sediments-and-fossils:fossil}{fossils}, same age (\omterm{def:g7:geological-time:fossil-succession}{principle of
fossil succession}). (b) The oldest is the sandstone with trilobites of
cliff A (it lies below the shale with ammonites, which is as old as the
bottom layer of cliff B). The youngest is the clay with leaves of cliff
B. (c) From the bottom up: sandstone with trilobites; shale with
ammonites; limestone with \omterm{def:g2:water-on-earth:ocean}{sea} urchins; clay with leaves.
\end{solution}

\begin{solution}{exo:g7:geological-time:11}
(a) The schist. It lies below the conglomerate (superposition), and the
\omterm{def:g3:stones-pebbles-sand:pebble}{pebbles} of the conglomerate are pieces of schist (inclusion); the
\omterm{def:g5:water-shapes-the-land:erosion}{erosion} surface also cuts the schist. (b) After: the sandstone was laid
down flat, so it was tilted later. (c) Before: the clay covers the \omterm{def:g7:earthquakes:fault}{fault}
without being broken. (d) The schist forms; it is worn away (wavy
surface); the conglomerate, the sandstone and the limestone are laid
down, in that order; the layers are tilted, and the \omterm{def:g7:earthquakes:fault}{fault} breaks them
--- the section alone does not say which of these two came first; the
tilted layers and the \omterm{def:g7:earthquakes:fault}{fault} are worn flat; the clay is laid down on top.
(e) Two: the wavy surface between the schist and the conglomerate, and
the surface between the tilted layers and the flat clay.
\end{solution}

\begin{solution}{exo:g7:geological-time:12}
(a) Half of 365 days is 182.5 days. January to June make
$31+28+31+30+31+30 = 181$ days, so the history is half over 1.5 days
into July: on \textbf{2~July at noon}. (b) $4540 \div 2 = 2270$ million
years. (c) $80 \div \num{4540000000} \times \num{31536000} \approx 0.56$:
about \textbf{0.6~second}.
\end{solution}

\begin{solution}{pb:g7:geological-time:1}
\textbf{1.} \num{4540} million years.

\textbf{2.} $4540 \div 365 \approx 12.4$ million years.

\textbf{3.} $\num{4540000000} \div (365 \times 24) = \num{4540000000}
\div 8760 \approx \num{518000}$ years (with the rounded \num{12.4}
million, $\num{12400000} \div 24 \approx \num{517000}$).

\textbf{4.} $\num{4540000000} \div \num{525600} \approx 8638$: about
\num{8600} years.

\textbf{5.} $4540 - 4030 = 510$ million years after the Earth formed;
$510 \div 12.4 \approx 41$ days. January has 31 days, so 41 days have
passed at the start of \textbf{11~February}.

\textbf{6.} $4540 - 3500 = 1040$ million years; $1040 \div 12.4 \approx
83.9$ (83.6 with the unrounded $12.44$): about 84 days. January and February
make 59 days, so it is about 25 days into March: \textbf{25~March}.

\textbf{7.} $66 \div 12.4 \approx 5.3$ days before the end of the year,
that is $365 - 5.3 = 359.7$ days after 1~January. The first eleven
months make 334 days: 25.7 days into December, on \textbf{26~December}.

\textbf{8.} $538.8 \div 4540 \approx 0.119$: \textbf{11.9~\%} of the
year is left, that is $0.119 \times 365 \approx 43$ days. $365 - 43 =
322$ days after 1~January: it begins in \textbf{November} (on
18~November).

\textbf{9.} $\num{10000} \div 8600 \approx 1.2$ minutes.

\textbf{10.} $5000 \div 8600 \approx 0.58$ minute, and $0.58 \times 60
\approx 35$: about \textbf{35 seconds} before midnight.

\textbf{11.} $\num{300000} \div 8600 \approx 34.9$: about \textbf{35
minutes} before midnight.

\textbf{12.} 35 minutes before midnight on 31~December: our species
appears at about \textbf{23:25}, in the last hour of the year.
\end{solution}
